\section{Introduction}

A dynamical-zeta candidate can possess an exact Euler germ and an operator
determinant without yet possessing the analytic structure relevant to a
Hilbert--P\'olya program.  The missing data include a source-native
archimedean completion, a functional equation, and a coherent reflection
center.  Classical Selberg and Ruelle zeta functions acquire such data from
trace formulas or geometric spectral identities
\cite{Selberg1956,Ruelle1976}.  The present arithmetic H\'enon construction
does not have an analogous global trace formula.

The preceding H\'enon sequence retained a genuine chronological
Fourier--cubic kernel, descended it through a normalized Galois trace, and
obtained a normalized-semifinite regularized determinant.  Its second,
third, and fourth logarithmic moments were then identified with a
genus-four curve, a Fano threefold packet, and a fivefold packet,
respectively.  Those identifications moved the normalized Euler object
successively to \(\RePart s>1/3\), \(\RePart s>1/4\), and
\(\RePart s>1/5\).  The second moment was additionally resummed into
modular elliptic factors.  This paper asks the next large question:
\emph{do the three cohomological moments admit one standard completion
without changing their prime clock?}

The answer has a positive and a negative part.  The positive structure is a
uniform pair of pure-weight packets.  If \(S_n\) is the Fermat cubic and
\(X_n\) is the source-ordered \((2,3)\) complete intersection, then
\[
 E_n=\Ql(0)\oplus H^{2n-2}_{\mathrm{prim}}(S_n)(n-1),
 \qquad
 O_n=H^{2n-3}(X_n)(n-2)
\]
have weights zero and one, and the exact moment trace is their sum.
Their total rank is \(4^n-1\).  This identity is a smooth-family theorem;
the H\'enon source is used only for \(n=2,3,4\), where smoothness away
from finite bad sets is inherited.

The negative structure comes from the clock rather than from a failure of
purity.  Field-degree normalization contributes
\[
 \frac1{p-1}=\sum_{j\ge1}p^{-j},
\]
so a standard factor is evaluated at \(u=ns+j\).  A pure weight-\(w\)
reflection consequently has center
\[
 s_{n,j}(w)=\frac{(w+1)/2-j}{n}.
\]
The leading odd factors align at zero, but the leading even factors do not,
and later denominator terms split the odd factors.  The exact
second-moment factorization already proves this bifurcation with factors
whose functional equations are known.  Thus the obstruction does not rely
on conjectural analytic continuation in the higher moments.

Three qualifications are essential.  First, center zero is not itself
incompatible with an RH normalization: recentering the completed Riemann
\(\xi\)-function produces a reflection about zero \cite{DLMF}.  Our
obstruction is the \emph{mismatch} of factorwise centers under one frozen
clock.  Second, the powers \(2/3\) and \(1/2\) are canonical
\(\Logzero\)-germs in a nonvanishing Euler domain, not meromorphic roots.
Third, the finite-rank no-go concerns a direct source-native
\(K\)-compatible packet preserving the same split-prime trace.  It does
not classify restriction of scalars, Galois-orbit counterpackets, or the
normalized-semifinite determinant already constructed.

The standard ingredients of the proof are classical.  We use Deligne
purity \cite{Deligne1974}, weak Lefschetz \cite{SGA2}, the
hypersurface Jacobian-ring description \cite{Griffiths1969II}, and
Hirzebruch--Riemann--Roch \cite{Hirzebruch1995}.  Serre and Deligne provide
the motivic local-factor, Tate-shift, and expected-completion conventions
\cite{Serre1970,Deligne1979}.  Buzzard--Gee and the automorphic
normalization literature explain why half-sum shifts are legitimate in a
different normalization category \cite{BuzzardGee2010,SST2014}; they do
not supply an ordinary half Tate motive preserving our coefficients.
Regularized-product approaches to Gamma and local factors
\cite{Quine1993,Deninger1991,Deninger1992} are close conceptual neighbors,
but no full H\'enon archimedean determinant is constructed here.

Our new contribution is the source-locked synthesis of these ingredients:
the exact two-weight trace and rank identity, the forced exponent \(2/n\),
the complete center tower, a scoped direct-compatible-system obstruction,
and an explicit Hodge projector gate.  Standard product stability in
motivic \(\lambda\)-rings \cite{Heinloth2007} explains why
denominator-cleared products are natural, but it does not prove the
functional equations of the present \(n=3,4\) packets.

The paper is organized as follows.  \Cref{sec:source} states the source and
main theorem.  \Cref{sec:weights} proves the point-count and rank identities.
\Cref{sec:logl} extracts the leading cohomological logarithm.
\Cref{sec:center} proves the weight--clock bifurcation and twist invariance.
\Cref{sec:compatible} gives the rank obstruction and the surviving odd
skeleton.  \Cref{sec:hodge} computes the fivefold Hodge ledger, and
\Cref{sec:route} records the Route-A status and claim boundary.
